\documentclass[12pt]{article}

% Packages
\usepackage[utf8]{inputenc}
\usepackage[T1]{fontenc}
\usepackage{amsmath}
\usepackage{graphicx}
\usepackage{natbib}
\usepackage{setspace}
\usepackage[margin=1in]{geometry}
\usepackage{hyperref}

% Double spacing
\doublespacing

% Title
\title{Hedging with Talent: How Selection Under Uncertainty Reshapes Automation's Workforce Effects}
\author{Matt Beane\\UC Santa Barbara \and Dan Sholler\\UC Santa Barbara}
\date{}

\begin{document}

\maketitle

\begin{abstract}
What happens to workers when automation arrives? Job insecurity theory predicts displacement: workers perceive threats and flee. Yet at a fulfillment facility deploying robotic automation, median tenure among temporary workers quadrupled---a reversal, not a reduction. Through a mixed-method study combining behavioral data on 59,021 temp worker separations across 21 automation deployments with 351 interviews and 718 hours of observation, we trace the organizational practices that produced this counterintuitive pattern. We find that novel technology created a specific epistemic condition: managers could not evaluate skill fit because the relevant performance criteria did not yet exist. Unable to select on skill, they defaulted to selecting on reliability---observable dependability---while deliberately shielding workers from automation's displacement implications. This \textit{reliability hedging} produced a compositional shift: reliable, low-turnover workers were channeled into automation-adjacent roles, generating tenure patterns that reversed standard predictions. These findings reveal that automation exposure is partly endogenous to organizational selection under uncertainty. Who encounters automation, and under what conditions, is not given by technology---it is produced by organizational practice.
\end{abstract}

\newpage

\section{Introduction}

Workers at a fulfillment facility stayed longer after the robots arrived. Not a little longer---dramatically longer. Median tenure among temporary workers increased from 19 hours to 86 hours, a gain of 67 hours representing more than a quadrupling of typical tenure. At a comparison facility introducing mature automation technology, tenure moved in the expected direction, dropping from 31 to 23 hours. The 74-hour differential between these patterns is large, robust, and directionally opposite to what theory predicts.

This finding is puzzling because the dominant narrative around automation and labor is one of displacement. Popular accounts warn of robots taking jobs \citep{frey2017future}. Academic research formalizes this expectation through job insecurity theory, which predicts that perceived threats to employment increase turnover intention and reduce organizational commitment, leading to accelerated separation \citep{sverke2002job, greenhalgh1984job}. The logic is straightforward: workers who perceive that their jobs are threatened should seek alternative employment, and those with viable alternatives should exit preemptively rather than wait to be pushed out. This prediction is especially strong for temporary workers, who face minimal barriers to exit. Temps can simply not return for their next shift, with negligible financial or reputational consequences. If any workers should flee impending automation, it is temps with their near-zero switching costs and minimal attachment to any particular employer.

Yet at the robot facility, temporary workers stayed. They did not merely fail to flee---they stayed dramatically longer than before the robots arrived, and dramatically longer than workers at a comparable facility deploying mature technology. The standard explanations for increased retention---improved job quality, stronger organizational commitment, reduced labor market alternatives \citep{mitchell2001turnover, griffeth2000meta}---find no support in the qualitative evidence. The robot facility did not become a more attractive workplace. Workers did not report greater satisfaction or describe robot work as especially engaging. Something else was operating.

The job insecurity literature cannot account for this pattern because it focuses exclusively on how workers respond to perceived threats, leaving unexamined the organizational practices that determine who encounters automation in the first place \citep{sverke2002job, greenhalgh1984job}. Labor market research recognizes that employment outcomes reflect bilateral sorting \citep{jovanovic1979job, autor2009studies}, but assumes that information about worker-task fit either exists or can be generated---an assumption that may fail when technology is genuinely novel \citep{barley1986technology, orr1996talking}. This suggests looking upstream: not just at how workers respond to automation, but at how they come to encounter it.

This paper traces the organizational practices that produced the counterintuitive tenure pattern. Through mixed-method analysis, we identify a mechanism invisible to the job insecurity framework---one operating before workers can perceive threats and respond. The account reveals how organizational selection under uncertainty shapes which workers encounter automation, and under what conditions \citep{orlikowski2000using}. The theoretical contribution emerges from this empirical tracing.

\section{Prior Work and the Displacement Puzzle}

\subsection{Job Insecurity and the Displacement Prediction}

The relationship between perceived job threats and worker behavior is among the most robust findings in organizational research. \citet{greenhalgh1984job} defined job insecurity as perceived powerlessness to maintain desired continuity in a threatened job situation, and subsequent research has consistently shown that this perception triggers withdrawal, disengagement, and ultimately separation. \citet{sverke2002job}'s comprehensive meta-analysis found a correlation of $\rho = .29$ between job insecurity and turnover intention, indicating that workers who perceive threats to their employment are substantially more likely to leave. Moderating factors---procedural justice \citep{brockner1994interactive}, threat salience \citep{ashford1989content}, communication quality \citep{dekker2002managing}---affect the strength of this relationship but not its direction. The fundamental causal model is consistent: threat perception drives exit behavior \citep{griffeth2000meta, hom2017review, mobley1979review}.

Applied to automation, this literature generates a clear displacement prediction: workers facing technological threats should exit faster than comparable workers not facing such threats. \citet{frey2017future}'s widely-cited analysis suggested that 47 percent of U.S. jobs were susceptible to computerization. \citet{acemoglu2020robots} documented that each additional robot per thousand workers reduced employment by approximately 0.2 percent in U.S. labor markets. Even \citet{autor2015jobs}'s more nuanced account, which noted that automation has historically complemented labor in aggregate, acknowledged that specific workers facing automation of their particular tasks face genuine displacement risk. The cultural and empirical narrative converge: automation threatens jobs, and threatened workers leave.

The displacement prediction is especially strong for temporary workers. Temps face minimal barriers to exit---they can simply not return for their next shift, with negligible financial or reputational consequences \citep{kalleberg2000nonstandard, kalleberg2009precarious}. They lack the organizational tenure, social ties, and job embeddedness that predict retention among permanent workers \citep{mitchell2001turnover, lee2004unfolding}. Employers use temporary staffing precisely because it provides flexibility \citep{houseman2001contingent}; from the worker's perspective, this same flexibility enables easy exit. If any workers should flee impending automation, it is temporary workers with their near-zero switching costs.

This combination of robust empirical findings and strong cultural narratives suggests that temporary workers confronting automation should flee. The prediction is theoretically grounded, empirically supported, and culturally reinforced. The prevailing expectation was unambiguous.

\subsection{What Prior Work Assumes}

The job insecurity framework rests on a specific causal model: automation creates objective threat, workers perceive threat, and perception drives exit behavior. The entire mechanism operates through worker cognition and response. This framing has dominated both academic research and popular discourse, producing a literature focused almost entirely on how workers interpret and react to technological change.

This focus on worker perception leaves organizational responses to technological uncertainty largely unexamined. Prior work assumes that workers encounter automation and then respond to it. The question is what shapes their response---factors like communication quality, procedural justice, and perceived control that affect how workers interpret threats \citep{brockner1994interactive, ashford1989content}. But the question of how workers come to encounter automation in the first place---how organizations allocate workers to automation-adjacent roles---has received far less attention.

This gap matters because automation deployment involves organizational choice. Someone must decide which workers will work with new technology. In contexts with labor fungibility---where workers can be assigned across roles---managers have discretion over these assignments. The exercise of that discretion could itself shape who experiences automation and under what circumstances. If selection into automation roles is systematic rather than random, the resulting patterns of who stays and who leaves may reflect selection as much as response.

The assumption that workers simply encounter automation and respond overlooks the organizational machinery that mediates that encounter. In large organizations with differentiated roles, workers do not automatically confront every new technology. Someone assigns them to positions. Someone decides which workers work in which areas. These decisions are not random. Managers make them based on some criteria---criteria that could themselves shape who experiences automation.

Prior work on labor market sorting provides relevant theoretical resources but proves ultimately inadequate for the phenomenon we document. Matching models \citep{jovanovic1979job}, temporary help agency research \citep{autor2009studies}, and personnel economics frameworks \citep{lazear1998personnel} all recognize that employment outcomes reflect bilateral processes of selection and sorting. But each assumes that information about worker-task fit either exists or will be generated through the employment relationship---the information problem is one of \textit{acquisition}, where fit information is costly but obtainable. The condition we encountered in the field was one of \textit{absence}: when technology is genuinely novel, the quality dimension itself is undefined. There is no ``match'' to reveal because the relevant performance criteria do not yet exist. The selection practice we document operates in this gap---before information about worker-technology fit can exist.

What happens when information is also asymmetric? When workers lack information about which roles involve automation and what automation might mean for their jobs, they cannot sort themselves based on automation preferences. Selection would then operate primarily through organizational discretion rather than worker choice.

Information asymmetry is especially likely in temporary work arrangements. Temps lack the organizational tenure that provides exposure to institutional dynamics. They lack the social networks that transmit informal information about what is happening and what it means. They lack the context to interpret signals about technology deployment. \citet{ashford2007guests} documented how temporary workers face ``liminality''---an in-between status that limits their access to organizational information. \citet{feldman1996underemployment} showed that temps often lack understanding of organizational dynamics that permanent employees take for granted. If information about automation implications does not reach temporary workers, they cannot self-select based on those implications.

The role of information asymmetry in labor allocation connects to broader research on how organizations manage uncertainty. \citet{galbraith1973designing} argued that organizations develop information-processing capacity to cope with uncertainty, including mechanisms for controlling what information reaches different organizational members. \citet{thompson1967organizations} noted that organizations buffer their technical cores from environmental uncertainty, which can include buffering workers from information about technological change. Applied to automation, these perspectives suggest that information about displacement implications might be deliberately or incidentally withheld from workers, affecting their ability to self-select.

\subsection{Theoretical Resources: Selection Under Uncertainty}

The sections above establish that job insecurity theory predicts displacement, and that prior work has focused on worker responses while leaving organizational selection largely unexamined. But a deeper question remains: when organizations \textit{do} select workers for novel technology roles, what determines the criteria they use? Existing theory offers resources for thinking about selection under uncertainty---resources that proved essential for interpreting the pattern we observed.

\textbf{Novel technology may eliminate skill-based selection.} When technology is genuinely novel, standard selection criteria may become unavailable. \citet{barley1986technology} demonstrated that identical technologies produce different role structures in different organizations because roles are constituted through practice, not determined by technology's properties. When novel technology arrives, the roles that will crystallize around it---troubleshooter, operator, coordinator---are unknowable in advance because they have not yet been enacted. \citet{orr1996talking} showed that even for established technology, the real knowledge required for competent work is invisible to management, embedded in the tacit practices and narratives of communities of practice \citep{lave1991situated}. For novel technology, this tacit knowledge has not been produced yet---the community of practice that would generate it does not exist. This raises a question: if skill-based selection becomes impossible, what criteria remain available?

This condition is distinct from uncertainty about which workers will succeed in known roles. It is closer to what \citet{tushman1986technological} called competence-destroying change, applied to personnel selection: the competence destroyed is the organization's ability to evaluate worker-task fit. The question becomes what managers do when matching cannot occur---not because information is costly to obtain, but because it has not been produced and cannot be produced prior to implementation.

\textbf{Temporary labor structure restricts available information.} The information architecture of temporary employment may further constrain selection options. Temporary work arrangements systematically generate data on limited dimensions of worker performance. Short expected tenure eliminates longitudinal observation---learning speed, adaptability, and flexibility all require repeated observation of the same worker over time, and temps do not stay long enough for such assessment \citep{bidwell2011paying}. Designed fungibility---the organizational construction of temp workers as interchangeable units---suppresses individual-difference data \citep{pfeffer1988taking, kalleberg2000nonstandard}. And agency mediation filters available information to attendance rates, no-show frequencies, and assignment completion rates, because these are the signals that temp employment structure allows agencies to observe \citep{autor2001wiring}. If technological novelty removes skill-based criteria and temp labor structure restricts available information, what criteria remain for selection?

One possibility: reliability---observable dependability---might be the only criterion that survives both constraints. Attendance and follow-through are context-independent; whether someone shows up consistently does not depend on what the work is. If this is so, then novel technology deployments staffed with temporary workers might show distinctive selection patterns.

\textbf{Satisficing and buffering as decision logics.} A third theoretical resource concerns what decision-makers do when facing constrained information. Facing novel problems without standard operating procedures \citep{cyert1963behavioral}, supervisors may search sequentially through available criteria and stop at the first that clears an acceptability threshold---the satisficing process that \citet{simon1947administrative} identified as the dominant mode of organizational decision-making. Reliability-based selection might also serve an organizational function: a novel process already operates at its uncertainty budget---task sequences are provisional, coordination patterns are unstable, error-recovery protocols do not exist. Selecting reliable workers could insulate the fragile technical core from additional disruption \citep{thompson1967organizations}.

These theoretical resources---epistemic constraints on selection, information architecture of temporary labor, satisficing and buffering---provide a framework for interpreting organizational responses to technological novelty. Whether and how these dynamics operate is an empirical question. The findings that follow document what we observed at facilities deploying novel versus mature automation technology.

\subsection{The Puzzle: Displacement in Reverse}

At the fulfillment organization we studied, the displacement prediction failed spectacularly. When robotic picking technology arrived, temporary workers stayed longer, not shorter. Median tenure increased from 19 hours before automation to 86 hours after---a gain of 67 hours, representing more than a quadrupling of typical tenure. This is not a null result suggesting that automation had no detectable association with tenure. It is a substantial reversal of the expected direction.

The magnitude of this reversal warrants emphasis. A quadrupling of median tenure represents a fundamental change in employment dynamics, not a marginal adjustment. Before automation, the typical temporary worker at Gallifrey stayed long enough to complete perhaps two or three shifts. After automation, the typical worker stayed long enough to complete training, develop familiarity with operations, and establish themselves as a known quantity to supervisors. These are qualitatively different employment relationships, not minor variations in duration.

The reversal becomes more striking when compared to a facility deploying mature automation technology. At a sorter facility, tenure moved in the expected direction, dropping from 31 hours to 23 hours after automation. The 74-hour differential between these facilities suggests that technology type matters: novel technology was associated with dramatically different tenure dynamics than mature technology. (Reported medians are rounded to whole hours; the 74-hour differential reflects exact computation from unrounded separation data.)

This comparison provides leverage for understanding the mechanism. If the tenure increase at the robot facility reflected general trends---improving labor markets, changing worker preferences, organizational improvements unrelated to automation---we would expect similar patterns at other facilities. We do not observe this. Tenure at the sorter facility decreased. The divergent patterns point toward something specific to novel technology, not general background factors.

This pattern cannot be explained by standard job insecurity theory. Both facilities introduced automation capable of performing tasks previously done by workers. Both represented potential displacement threats. If anything, the robot facility should have triggered stronger flight responses---robots are more culturally salient as displacement threats than sorting equipment, and the robot deployment received more attention within the organization. The cultural salience of robots makes the reversal more puzzling: workers facing the more threatening-appearing technology stayed longer, while workers facing less threatening-appearing technology fled.

The puzzle deepens when we consider the counterfactual. If workers at the robot facility had perceived displacement threats and responded according to job insecurity theory, tenure should have decreased. If workers had been indifferent to the robots, tenure should have remained stable. Instead, tenure quadrupled. Something actively increased retention---not by a small margin, but dramatically. The pattern requires explanation, not merely acknowledgment.

Consider three possible interpretations. First, workers may have found robot work intrinsically appealing, preferring it to traditional warehouse work and therefore choosing to stay. Second, working conditions may have improved in ways that increased retention independently of automation. Third, the composition of workers in automation-adjacent roles may have shifted toward workers predisposed to stay. The first two interpretations focus on worker responses to changed conditions. The third focuses on organizational selection. The findings that follow provide evidence most consistent with the third interpretation.

The puzzle is not whether job insecurity theory is wrong. The sorter facility results suggest it applies in at least some automation contexts. Workers at Skaro, facing mature technology, showed the expected pattern: tenure decreased after automation deployment. The displacement prediction held for mature technology. The puzzle is why novel technology---robots---produced the opposite pattern. What is it about technological novelty that inverts displacement dynamics?

This paper traces the organizational practices that produced the observed pattern. The findings that follow document how managers responded to the uncertainty created by novel technology, and how those responses shaped which workers encountered automation under what conditions. The theoretical contribution emerges from this empirical account.

\section{Methods}

\subsection{Research Setting}

The study takes place at a third-party logistics company, which we call UNIT, operating fulfillment facilities across the United States. UNIT provides warehousing, inventory management, picking, packing, and shipping services for client companies ranging from small e-commerce operations to large retailers. At the time of the study, UNIT operated over twenty facilities nationally, ranging from small specialty operations handling specialized products to large distribution centers processing millions of units annually.

The fulfillment industry provides an ideal context for studying automation and labor dynamics. Fulfillment work is characterized by tasks that are increasingly susceptible to automation---picking items from inventory, sorting packages by destination, packing items for shipment---combined with workforce characteristics that create high turnover and labor market fluidity. The industry has been a leading adopter of warehouse automation, with companies deploying a range of technologies from traditional conveyors and sortation systems to cutting-edge autonomous mobile robots. This makes fulfillment a bellwether for understanding how automation affects workers more broadly.

The work itself is physically demanding and repetitive. Pickers walk miles per shift navigating warehouse aisles to retrieve items from shelving units. Packers stand at stations assembling boxes, applying labels, and ensuring items are properly protected for shipment. Sorters move packages from conveyor lines to staging areas, often lifting and carrying heavy loads throughout the shift. The physical demands contribute to high turnover, as do wages that typically hover near regional minimums for warehouse work. Shift schedules vary but commonly include early mornings, late nights, and weekend work, further contributing to workforce instability.

UNIT employs a mix of permanent employees and temporary workers. Permanent employees hold full-time positions with benefits and job security protections typical of stable employment relationships. They accumulate organizational tenure, develop relationships with supervisors and coworkers, and have access to information networks that transmit knowledge about organizational dynamics. Temporary workers are employed through staffing agencies, working at UNIT facilities on an assignment basis without the benefits and protections afforded permanent employees. Temps can be terminated at will, receive minimal benefits, and have no expectation of continued employment beyond immediate operational needs.

Temps comprise the majority of the frontline workforce at most facilities, typically between 85 and 95 percent depending on season and facility needs. This reliance on temporary labor reflects industry-wide patterns in fulfillment, where demand volatility and the physical demands of the work create high turnover that makes permanent staffing costly. Turnover among temps is extremely high---many workers separate within days or weeks of starting---creating a continuous churn of hiring, training, and separation. The staffing agencies that supply temps maintain large pools of potential workers, enabling rapid replacement when workers leave. From UNIT's perspective, the temp model provides operational flexibility. From the worker's perspective, it provides easy entry and exit---a job can be obtained quickly, and leaving requires simply not returning for the next shift.

Between 2019 and 2021, UNIT deployed automation technologies at multiple facilities. These deployments varied in technology type, creating the variation that enables the current study. Some facilities received robotic picking systems---autonomous mobile robots that transport inventory shelves to stationary workers, a relatively novel technology class that emerged commercially in the mid-2010s. These robots navigate warehouse floors using sensors and software, bringing products to workers rather than requiring workers to traverse the facility to retrieve products. Other facilities received automated sorting equipment---conveyor systems with optical scanning and diverters that route packages to appropriate destinations, a mature technology class with decades of industry deployment history. Sortation technology is well-understood; organizations know what skills predict success, and equipment suppliers provide training protocols and staffing guidance.

The primary comparison involves two facilities: Gallifrey, which received robotic picking technology in early 2020, and Skaro, which received automated sorting equipment in early 2020. These facilities are similar in size, workforce composition, and operational focus, differing primarily in the type of automation deployed. Both serve regional distribution needs, employ several hundred workers during peak periods, and process comparable volumes of units. The key difference is technology type: novel robots at Gallifrey versus mature sorters at Skaro. Additional facilities provide robustness checks, allowing examination of whether patterns observed at the primary comparison sites generalize across the organization.

\subsection{Data}

The analysis draws on three primary data sources combining quantitative administrative records with qualitative interview and observation data.

The first data source comprises administrative records on temporary worker separations. These records cover 59,021 temp worker separations from 2017 through 2024 across all UNIT facilities. Each record includes the worker's hire date, separation date, separation type, facility, and total hours worked. Hours worked provides the primary measure of tenure---the length of time a worker remained employed before separation. The administrative data were extracted from UNIT's human resources information system, which tracks all employment events. The long time series enables examination of tenure patterns before, during, and after automation deployments, providing both baseline comparisons and post-deployment trajectories.

The second data source comprises administrative records on permanent employee terminations. These records cover 2,506 permanent employee separations from 2019 through 2022 (permanent-employee records were not systematically archived prior to 2019). Comparing tenure patterns across temporary and permanent workers allows examining whether the selection pattern operated differently for temps versus permanents. Permanent employees have characteristics---longer tenure, greater information access, reduced labor fungibility---that should attenuate the selection mechanism. If permanent employees showed the opposite pattern from temps, it would provide evidence consistent with the proposed mechanism.

The third data source comprises qualitative data from an extended field study. Between 2019 and 2021, we conducted 718 hours of observation across 54 site visits at 12 warehouse facilities across six warehousing firms. Four of these facilities belonged to UNIT; the others provided comparative context for understanding industry practices. During this ethnographic work we performed 351 semi-structured interviews with workers, supervisors, managers, and executives. Interviews were conducted onsite during facility visits and remotely via video call. Interview protocols covered work processes, reactions to automation, selection and assignment practices, information flows, and general organizational dynamics. Interview duration ranged from 15 minutes for brief frontline worker conversations to 90 minutes for in-depth management interviews, with most lasting 30-45 minutes.

Interview protocols evolved as the study progressed, following recommended practice for inductive qualitative research \citep{eisenhardt1989building}. Early interviews used broad questions about work experience and organizational change: How did you come to work here? What is a typical day like? What has changed recently? As the selection mechanism emerged as a potential explanation for the quantitative patterns, we added targeted questions about assignment practices, selection criteria, and information flows: How are workers assigned to different areas? What do you look for when selecting workers for new roles? What did workers know about the automation project? This iterative refinement enabled us to probe the mechanism while remaining open to alternative explanations.

Observations focused on work practices, interaction patterns, and physical arrangements of work. We observed shift changes, training sessions, team meetings, and routine operations. At facilities with automation, we spent concentrated time in both automated and non-automated areas, comparing workflow and worker behavior across contexts. We observed managers making assignment decisions, watched how information about projects was communicated, and noted which workers were placed in which areas. Field notes were typed within 24 hours of each site visit, following \citet{miles2014qualitative}'s recommendations for maintaining data quality. Notes included both descriptive content and analytical memos capturing emerging interpretations.

The qualitative data serve multiple purposes in this study. They provide mechanism evidence---direct observation of the selection criteria managers used when assigning workers to automation roles. They document information practices---how managers communicated about automation and what information reached workers. They capture worker perspectives---how workers understood their roles and whether they were aware of automation implications. Together these purposes allow triangulating the quantitative tenure patterns with qualitative evidence about the processes producing those patterns.

This mixed-method design follows \citet{edmondson2007methodological}'s guidance on methodological fit: the nascent theory state of our research question called for inductive qualitative work to identify mechanisms, combined with quantitative data to establish the pattern requiring explanation. The combination is essential. Quantitative data alone would reveal the tenure increase without explaining it. Qualitative data alone would document organizational practices without establishing whether those practices were associated with distinct workforce outcomes. Together they enable both pattern establishment and mechanism identification.

\subsection{Analytic Approach}

The quantitative analysis examines tenure distributions before and after automation deployment at facilities with different technology types. The primary comparison uses difference-in-differences logic: we compare the change in tenure at the robot facility to the change in tenure at the sorter facility, isolating the effect of technology type from general temporal trends. If both facilities showed similar tenure changes, we could not attribute the pattern to technology type. The divergent patterns---substantial increase at the robot facility, decrease at the sorter facility---suggest technology type matters.

Tenure distributions in temporary labor are highly right-skewed, with many workers separating after minimal hours and a long tail of workers with extended tenure. The distribution is not approximately normal; means are pulled upward by long-tenured outliers. Median tenure provides a more robust summary statistic than mean tenure in this context. We report median tenure changes as the primary quantitative results. Statistical significance is assessed using Mann-Whitney tests for comparing distributions across periods, which do not assume normality and are appropriate for skewed distributions.

The qualitative analysis follows abductive principles, iterating between data and theory to develop explanations for observed patterns \citep{locke2007grounded}. Initial coding identified passages related to selection criteria, automation awareness, and worker responses. We coded for who made selection decisions, what criteria they reported using, how workers described their assignment processes, and what workers knew about automation implications. Focused coding then examined how selection operated---what criteria managers used, how assignments were made, what workers knew---with particular attention to whether the evidence supported or contradicted the emerging selection interpretation.

The mixed-method design serves an explanatory purpose. The quantitative analysis establishes the pattern requiring explanation---tenure increases at novel technology facilities but not mature technology facilities. The qualitative analysis provides explanation---the selection mechanism that produces this pattern. Each method provides what the other cannot: quantitative breadth and rigor combined with qualitative depth and mechanism visibility.

A limitation of this design warrants acknowledgment at the outset. The primary quantitative comparison rests on a single novel-technology site---Gallifrey is the only facility in the data that deployed robotic picking technology during the observation window. We interpret the converging quantitative patterns and qualitative mechanism evidence as identifying a theoretically grounded selection process rather than establishing causal identification. The mechanism points to specific conditions under which the pattern should and should not appear, which we articulate in the Discussion.

\section{Findings}

The findings that follow trace how managers at a robot facility responded to technological uncertainty, and how that response produced the counterintuitive tenure patterns observed in the administrative data. The account proceeds in three stages. First, we document what we call \textit{reliability hedging}---the practice of selecting workers for novel technology roles based on observable dependability rather than skill fit. Second, we examine \textit{information shielding}---the deliberate withholding of automation implications that enabled managerial selection to operate without triggering worker flight. Third, we show how these practices produced a \textit{compositional shift}: the tenure increase reflected not workers choosing to stay, but reliable workers being channeled into automation-adjacent roles. The quantitative patterns become interpretable only after understanding the organizational practices that produced them.

\subsection{Reliability Hedging}

When I0 learned he would lead the robot deployment team at Gallifrey, he faced a problem with no obvious solution. Senior management had told him how many workers he needed---pickers, line captains, drivers---but not what kind of workers. The technology was genuinely new. No one at UNIT had worked with autonomous mobile robots. There was no institutional knowledge about what skills or attributes predicted success with the system. I0 was building a team for work that had never been done before.

His recounting of receiving the assignment captures this uncertainty:

\begin{quote}
So N01 said, `Hey, I0. You're going to be in charge of forming this team. These are the people that we need.' I believe it was either yourself or maybe [inaudible] gave us a list of how many people that they foreseen was needed for the project. So I got issued, `You need these many pickers. You need these many line captains. You need these many drivers, yada, yada, yada.' So I went out. I get my chart. And the first thing I did is I verified attendance across first and second shift, but it's probably the most important thing, right? You need to show up to work.
\end{quote}

The directive specified quantities but not qualities. And I0's immediate response---before anything else---was to check attendance records. Not technical aptitude. Not interest in technology. Not problem-solving ability. Attendance. The selection criterion for working with novel automation was whether someone could be relied upon to appear.

This choice was not arbitrary. It reflected a specific logic about how to select under epistemic impossibility. I0 explained his reasoning directly when asked why he prioritized reliability:

\begin{quote}
At the end of the day, I didn't want to go out there and start shooting off names, picking up people because I'll be honest with you, I don't know if the guys going to be the best guy or if he's going to be the worst guy. T4 was able to answer that question and she did it very effectively. I said, `I need you to help me draft a team of the best possible people that one, they're going to be here on time, and they know these individual areas.'
\end{quote}

I0's statement reveals the core dynamic: he could not identify who would be ``best'' or ``worst'' for robot work because the criteria for success were unknown. The technology was too new. Nobody had experience. In the absence of skill-based selection criteria, he defaulted to the most observable signal of general competence available in temporary labor---whether someone showed up consistently and followed through on commitments.

We term this practice \textit{reliability hedging}. Like financial hedging, which reduces exposure to uncertain outcomes by diversifying across assets likely to perform adequately under multiple scenarios, reliability hedging reduces exposure to uncertain skill requirements by selecting workers likely to perform adequately regardless of what the technology demands. A worker who reliably appears for shifts will be present to learn whatever the new technology requires. A worker who completes assignments without quitting mid-engagement will stay long enough to develop competence. Reliability serves as a hedge against the unknown.

The practice extended beyond I0's initial selection. To identify specific candidates, he sought help from T4, a supervisor with ground-level knowledge of the workforce. ``So I went over there and got with T4,'' I0 recalled. ``Because by the way, she was one of the ones concerned about the robots.'' Even someone worried about automation's implications was enlisted to help identify reliable workers for automation roles. The selection criterion transcended attitudes toward the technology itself.

T4's involvement reveals how reliability assessment drew on multiple sources of knowledge. As a line supervisor, she had day-to-day visibility into worker behavior that I0, operating at a higher organizational level, lacked. She knew who showed up consistently, who completed tasks without supervision, who could be counted on during peak periods. Her local knowledge enabled the translation of I0's reliability criterion into specific worker selections. The process combined hierarchical direction (reliability as the selection criterion) with ground-level knowledge (who actually met that criterion).

N01, a supervisor at Gallifrey who managed one of the robot lines, independently described the same constraint. He had tried to build stable teams by assigning the same workers to the same lines---but attendance made this impossible:

\begin{quote}
What I had implemented at first was having the same people on each line, but attendance became a problem because under the [staffing agency] model, they controlled basically who went to what line. People wasn't here every day.
\end{quote}

N01's account reveals that reliability was not merely I0's personal preference but a structural response to the labor context. The staffing agency model generated unreliable attendance, and when supervisors tried to form coherent teams for novel work, attendance was the binding constraint. ``You just don't know what to get,'' N01 observed about agency workers assigned to robot lines.

Other supervisors described similar approaches. C19, a team lead at a different facility, articulated the logic with particular clarity:

\begin{quote}
You've got to have people that you can rely on. I'm not here 24 hours a day, so on second and third shift you're by yourself\ldots I need adults.
\end{quote}

When asked what he looked for in workers, C19 emphasized tenure stability: ``Someone that's not constantly job jumping. That's maybe just me. But it's kind of a red flag, one of my red flags.'' Another supervisor characterized the decision logic succinctly: ``Tech skills you can train. Showing up, you can't. If someone's reliable and they don't know the system yet, you work with that. If someone's unreliable, doesn't matter how smart they are---they're not going to be there when you need them.''

This perspective inverts the usual priority. Technical capability is typically valued because it enables immediate productivity. But under systematic legibility collapse---when what capabilities matter is unknowable---reliability becomes paramount because it enables learning and adaptation. A reliable worker without robot skills can acquire those skills through training and experience. An unreliable worker, regardless of aptitude, may not be present when needed.

The contrast with mature technology deployment illuminates the distinctiveness of reliability hedging. At Skaro, where sortation equipment was deployed, managers described different selection approaches. The technology was well-understood. Industry experience provided templates for what skills predicted success. A supervisor at Skaro explained the difference:

\begin{quote}
With the sorters, we know exactly what we need. You need people who can keep pace with the line, who can handle the scanning equipment, who don't get flustered when packages are coming fast. We've been running sorters for years. There's a profile that works. With something brand new, you're guessing.
\end{quote}

This statement captures the difference between mature and novel technology selection. Mature technology allows skill-based selection because the relevant skills are known. Novel technology forecloses this option, pushing managers toward reliability as a default criterion.

An important qualification: selection was not purely single-criterion. While reliability served as the primary screen---the necessary condition that workers had to clear before other criteria mattered---supervisors also considered secondary attributes among the reliability-screened pool. A senior manager at a sorter facility described looking for ``eagerness to go and do different things, and ones that we knew had an aptitude for learning new things quickly,'' though he immediately added: ``Really, it boils down to effort, though, to be honest with you.'' Another manager identified curiosity as the distinguishing factor among reliable workers: ``You want somebody with a good work ethic, but the thing that's really come to light more recently is somebody who has curiosity\ldots they want to know why things are done.'' He estimated that roughly 20 percent of workers were ``driven and curious,'' while 50 percent simply ``come to work every day''---the reliable baseline. I0's own selection statement hints at this multi-criterion logic: ``the best possible people that one, they're going to be here on time, \textit{and they know these individual areas}''---the second criterion is competence, not reliability.

This multi-criterion structure matters for the mechanism. Reliability was the \textit{necessary} first filter because it was the only criterion that survived both the legibility collapse created by novel technology and the information constraints of temporary labor. But once the reliability screen was passed, supervisors considered additional attributes---curiosity, engagement, area knowledge---that are \textit{not} definitionally linked to retention. The selection process was a portfolio rather than a single bet, which makes the hedging metaphor more apt: managers diversified across traits that might predict success under unknown conditions, with reliability as the anchor position.

The practice of reliability hedging has a critical implication for understanding tenure patterns. Workers selected for reliability are disproportionately workers who do not quit---consistent attendance and follow-through overlap substantially with the behaviors that predict retention. But the multi-criterion evidence partially dissolves what might otherwise be a tautological link. Workers selected on curiosity or eagerness to learn are not definitionally low-turnover; their retention reflects something beyond mechanical reliability-retention overlap. If the tenure increase were driven entirely by the tautological reliability-retention linkage, we would expect no evidence of secondary selection criteria---managers would simply pick whoever showed up most and the tenure pattern would follow mechanically. Instead, the qualitative evidence documents a richer selection process in which reliability opened the door but curiosity, engagement, and competence determined who walked through it. These secondary criteria introduce genuine variance in why selected workers stayed: a curious worker who remained in a robot role may have stayed because the work engaged her, not merely because she was dispositionally unlikely to quit. The compositional shift therefore reflects a combination of mechanical reliability-retention overlap and authentic worker-technology engagement discovered through multi-criterion selection---a combination that makes the tenure increase partly but not wholly tautological.

\subsection{Information Shielding}

Reliability hedging could only operate if workers were not simultaneously self-selecting based on automation preferences. If workers knew which roles were automation roles and what automation might mean for job security, different dynamics would emerge. Risk-averse workers might flee before being selected. Workers with technical ambitions might volunteer. Either response would contaminate the reliability-based selection with self-selection effects.

The qualitative evidence reveals that managers actively prevented this contamination through what we term \textit{information shielding}---the deliberate withholding of automation implications from workers. Information shielding was not incidental to reliability hedging; it was an enabling condition that preserved managerial control over the selection process.

The robots arrived at Gallifrey before management had formulated a communication plan. I0 described the resulting situation:

\begin{quote}
This robot got pulled out of the box, and we started doing basically the physical setup of the equipment. There was a lot of concern. So that kind of spread like wildfire within, I'd say, the first four hours that K1 was on site. So we went ahead and took the initiative, grouped up first and second shift and said, `Hey, guys, you know what? There's going to be some robotic setup and some testing on [inaudible]. Just be transparent as possible. This is not for removal of your position. This is to optimize our throughput and production.'
\end{quote}

The messaging acknowledged the robots' presence while explicitly denying displacement implications. Workers were told the robots were for optimization, not job elimination. The choice of framing was deliberate---``transparent as possible'' meant transparent about the robots' existence, not about their implications. This framing served to contain worker anxiety without revealing the labor cost reduction that motivated the technology investment.

When I0 reflected on what he had communicated versus what he understood to be true, the gap was stark:

\begin{quote}
And I just kind of left it open like that because, I mean, the ultimate goal here is to reduce labor costs and optimize our output. But I'm not going to go ahead and deliver that message. That is executive leadership's call if that's going to move forward like that.
\end{quote}

I0 understood that automation was ultimately about reducing labor costs---the classic displacement story that job insecurity theory describes. He chose not to communicate this understanding. The decision was explicit, deliberate, and recognized as such. Workers received a sanitized message about optimization while the displacement implications remained hidden. I0's statement also reveals the organizational politics of disclosure: he deferred to executive leadership on whether to reveal the labor cost implications, effectively passing responsibility upward while maintaining the information shield at his level.

The information asymmetry was pronounced. When asked what workers typically understood about automation projects, a senior manager at another facility characterized the pattern:

\begin{quote}
They know something's happening. They can see the equipment. But they don't connect the dots to headcount. That's not---we don't talk about that with floor associates. It's not really their concern until it affects them directly, and by then the decisions are already made.
\end{quote}

This statement reveals a managerial perspective that treats automation implications as information workers do not need. The logic positions displacement decisions as matters to be handled by management, with workers informed only when decisions directly affect them---at which point, as the manager noted, ``the decisions are already made.'' Information shielding thus preserves managerial control over both the automation implementation and the communication of its implications.

The effectiveness of information shielding was evident in worker interviews. When asked about the robot project, frontline workers consistently reported having minimal information. C7, a line leader, provided a representative response:

\begin{quote}
I have no idea. All I know is they said, `We're bringing robots in.' I don't know what role they're going to play or what associates will be dismissed because of the robots. I have no idea.
\end{quote}

C7's statement captures the information environment that information shielding created. Workers knew robots existed. They did not know what the robots meant for their jobs. They could not self-select based on automation implications because they did not know what those implications were. This information deficit enabled managerial selection to operate uncontaminated by worker preferences.

The worker experience of being selected into a robot role illustrates this opacity. O7, a temporary worker at another facility, had applied for a regular warehouse position when the staffing agency redirected her:

\begin{quote}
They had called me from the actual temp service instead of here and told me they had this unique position available for me that basically was just a startup type program and wanted to see if I was interested. And so that's exactly what---because I'm like, `What the heck? Let me try this.'
\end{quote}

O7 had no prior pick-and-pack experience and received a twenty-minute crash course from a trainer who was leaving the country that day. ``There was really no one that had done this type of position,'' she recalled. ``I had to see the ins and outs. And I kind of put it together of how it was going to work\ldots it was just learn as you go type stuff.'' From the selected worker's perspective, assignment to a robot role was not a deliberate career choice informed by understanding of what automation meant---it was an opaque organizational process that deposited her in a position she had never heard of, with training that barely existed. K13, a temp at a different facility, captured this perception of selection as arbitrary: ``I think they just pick out a random and just see how it goes.''

The pattern extended beyond Gallifrey. At a facility preparing for a later robot deployment, supervisors R10 and N29 described learning about their own facility's automation plans by accident:

\begin{quote}
We found out about three weeks [ago]\ldots By accident, by accident, they told us. It was not like ``Oh, by the way, we are changing [plans].'' And we're on this weekly calls, update on what the facility, no mentioned of it, right.
\end{quote}

Their colleague added: ``He literally didn't tell no one. It was just out of the blue.'' Even supervisors on weekly update calls were not informed about major robot deployments at their own facilities. This pattern suggests that information shielding operated through two distinct layers. The first is \textit{structural information poverty}: temporary workers---and even some supervisors---lack the organizational tenure, communication network access, and institutional context to interpret signals about technological change. I0 himself was uninformed about robot plans until one month before deployment. The structural layer is a scope condition that makes shielding possible; it generalizes across organizations that rely on temporary labor. The second layer is \textit{deliberate withholding}: I0's explicit choice not to ``deliver that message'' about labor cost implications. This deliberate layer is a practice that exploits the structural condition---a contingent managerial choice rather than a structural inevitability.

Other frontline workers echoed this information vacuum. A picker at Gallifrey, when asked what the robots would mean for her job, responded: ``Nobody's told me anything about that. I just come in and do my job.'' A packer at another facility similarly reported: ``They don't really explain things to us. New equipment shows up, they train you on it, that's it.'' The consistency of these responses across workers and facilities suggests that information shielding reflected both structural information poverty inherent in temporary labor arrangements and deliberate withholding practices that varied across managers but converged on the same outcome: workers who could not self-select based on automation preferences.

The combination of reliability hedging and information shielding produced a specific organizational dynamic. Managers selected workers for robot roles based on observable dependability. Workers did not know they were being selected for robot roles or what such roles might mean. The selection process operated through managerial discretion rather than bilateral sorting. Reliable workers were channeled into automation-adjacent positions without the opportunity to flee or volunteer based on automation preferences. The two layers of information shielding---structural and deliberate---ensured that the selection remained uncontaminated by worker self-selection.

\subsection{The Compositional Shift}

The practices of reliability hedging and information shielding suggest a specific interpretation of the tenure patterns in the administrative data. If managers selected reliable workers for robot roles, and reliable workers are disproportionately low-turnover workers, then the tenure increase at Gallifrey reflects a \textit{compositional shift} in who worked in automation areas rather than a change in how workers responded to automation exposure.

The quantitative patterns are consistent with this interpretation. At Gallifrey, median tenure increased from 19 hours before robot deployment to 86 hours after---a gain of 67 hours representing more than a quadrupling of typical tenure. At Skaro, the sorter facility deploying mature automation technology, tenure moved in the expected direction, declining from 31 to 23 hours. The 74-hour differential between these patterns is substantial.

The magnitude warrants emphasis. A 74-hour difference in median tenure is not a subtle effect requiring sophisticated detection. It represents fundamentally different turnover dynamics. At Gallifrey, the typical worker stayed long enough to complete meaningful training and establish presence in the facility. At Skaro, the typical worker separated before completing basic orientation. These are qualitatively different employment relationships, not marginal variations in duration.

The compositional shift interpretation gains support from within-facility variation at Gallifrey. Post-deployment, supervisors showed dramatically different tenure profiles for their workers. Supervisors overseeing what appear to be robot operations---identified by median tenure exceeding 300 hours---had workers with median tenure of 626 hours. Supervisors overseeing non-robot operations had workers with median tenure of 101 hours. Workers in robot areas stayed six times longer than workers elsewhere in the same facility.

This within-facility variation is difficult to explain through facility-level factors like management quality or local labor markets---both areas faced identical facility conditions. It is consistent with the compositional shift interpretation: robot areas received selected workers who were reliable and thus low-turnover, while traditional areas received the broader workforce with typical turnover patterns.

The shift in workforce composition appears directly in the tenure distribution. Before robot deployment, only 5 percent of separations at Gallifrey involved workers with 600 or more hours of tenure. After deployment, this proportion rose to 23 percent---nearly a fivefold increase in the share of long-stayers. The facility-wide median of 86 hours reflects the mixture of two distinct populations: a smaller group of robot-area workers with very high tenure and a larger group of traditional workers with shorter tenure. The 50th percentile of the pooled distribution shifted upward because of this compositional change. The tenure increase emerges from changing the mix, not changing individual behavior.

At Skaro, by contrast, the share of long-stayers barely changed (3.2 to 4.4 percent). Mature technology did not trigger reliability hedging because skill requirements for sorter equipment were understood. Managers could select based on relevant experience and aptitude rather than defaulting to reliability. Standard displacement dynamics prevailed.

Additional sorter facilities showed the same pattern. Across six facilities with sufficient pre/post data, median tenure changes ranged from -10 to +10 hours, averaging approximately +1 hour. No sorter facility showed the dramatic increases observed at Gallifrey. The pattern appears specific to novel technology.

The permanent employee pattern provides additional evidence. If the compositional shift operates through reliability hedging enabled by labor fungibility and information shielding, it should not operate for permanent employees who have less fungibility and more information access. Corroborating this interpretation, permanent employees showed the opposite pattern from temps at Gallifrey: their tenure decreased rather than increased after robot deployment.

The divergence makes sense under the compositional shift interpretation. Permanent employees had what temps lacked: information. With longer organizational tenure, they had witnessed prior rounds of technology change. With better access to communication networks, they could assess what robot deployment signaled for the facility's trajectory. Information shielding that kept temps uninformed was less effective for permanents. Their tenure decrease is consistent with standard displacement dynamics operating on workers who understood automation's implications, while the temp tenure increase reflects compositional change among workers who did not.

\subsection{Alternative Explanations and Robustness}

The selection mechanism we propose is not the only possible explanation for the observed tenure patterns, and several alternative explanations deserve consideration. The most significant threat is that Gallifrey differs from other facilities in ways unrelated to technology novelty---perhaps unusually effective management or a favorable local labor market. Several observations mitigate this concern: Gallifrey's pre-automation tenure was among the lowest in the organization (19 hours median), the within-facility variation between robot and non-robot areas is difficult to explain through facility-level confounds, and the permanent employee pattern showed the opposite direction from temps. The timing of deployment also matters: robots arrived on January 20, 2020, and the tenure increase appeared within the first ten days---well before COVID-19 lockdowns began in mid-March. The pre-COVID diff-in-diff was +115 hours, larger than the full-sample estimate, while tenure at Skaro decreased by 24 percent during the same period.

Treatment effects represent a second alternative: workers in robot roles may have found the work intrinsically more appealing and stayed because they liked it, not because they were selected for reliability. We cannot fully separate selection from treatment with observational data, but the magnitude (tenure quadrupling) seems too large for pure treatment effects given similar core work tasks, interview evidence suggests workers did not find robot work especially appealing, and permanent employees in robot areas did not show increased tenure. Some worker self-selection likely occurred despite information control---the mechanism probably operates through a combination of managerial selection and limited self-selection rather than pure managerial discretion, though this does not change the core insight that selection concentrates reliable workers in automation roles. Finally, regression to the mean is possible given Gallifrey's very low baseline, but the magnitude of the increase substantially exceeds what regression alone would predict, and Skaro's decrease is inconsistent with a pure regression explanation. The evidence is most consistent with the selection mechanism we propose.

\section{Discussion}

\subsection{Theoretical Contribution}

The central contribution of this paper is to identify an upstream mechanism that shapes which workers encounter automation before worker choice can operate. The tenure pattern at Gallifrey---temps staying four times longer after robots arrived---is evidence of this selection process, not the headline finding. The headline finding is the mechanism itself: when technological novelty eliminates skill-based selection criteria, and the information architecture of temporary labor restricts available criteria to reliability, managers satisfice on reliability as a criterion that also buffers fragile novel processes from compounding uncertainty. The resulting compositional shift---reliable workers channeled into automation-adjacent roles---produces tenure patterns that reverse standard displacement dynamics.

This contribution is different in kind from prior work on job insecurity, which has identified moderators affecting how workers respond to perceived threats \citep{brockner1994interactive, ashford1989content}. Those moderators operate within the standard causal model: automation creates threat, workers perceive threat, perception drives exit. The mechanism documented here intervenes before this sequence can unfold. It operates upstream, through organizational practices that shape who encounters automation rather than through worker perception and response. Identifying this upstream selection mechanism is the paper's theoretical move.

The three-layer theoretical architecture---legibility collapse, information architecture, satisficing-plus-buffering---adds precision to when and why selection effects should dominate treatment effects in automation contexts. The mechanism is specific to \textit{skill uncertainty}, not to initiative importance. The Skaro comparison demonstrates this: deploying mature automation technology was equally important organizationally but did not trigger reliability hedging because skill requirements were known. This specificity distinguishes the mechanism from a generic ``important initiatives trigger careful selection'' story and points to specific conditions under which the pattern should and should not appear. Where novel technology meets temporary labor, we would expect reliability-based selection to emerge; where the same technology meets permanent labor, the richer information environment should partially compensate; where mature technology meets either labor type, skill-based selection should prevail.

This reframing has implications for interpreting aggregate employment patterns around automation. When tenure increases or employment remains stable at firms deploying technology, the standard interpretation is that automation creates good jobs or that workers prefer automation-adjacent work. The compositional shift mechanism suggests an alternative: stable employment may reflect who is selected into automation roles rather than how workers respond to automation. If reliable workers---inherently less likely to quit---are systematically channeled into automation areas, aggregate tenure will increase regardless of whether workers find the work attractive.

The contribution also connects to broader debates about technology and work. \citet{orlikowski2000using} argued that technology effects emerge through the practices of those who implement and use technology. Reliability hedging and information shielding illustrate this in the labor allocation domain: identical automation technology could be implemented with random assignment, skill-based selection, or reliability-based selection, producing different tenure outcomes. The practices documented here are not properties of the technology but organizational responses to the epistemic condition that novelty creates. As organizations deploy increasingly novel technologies while relying on contingent workforces---gig work, platform labor, temporary staffing---the conditions for reliability-based selection may become more common. Understanding how technological novelty interacts with labor market structures to shape who encounters automation becomes increasingly important.

\subsection{Boundary Conditions}

The practices documented here operated under specific conditions that define their boundary. These conditions are not universal; they characterize specific organizational and labor market contexts. Understanding when the mechanism operates---and when it does not---is essential for assessing its broader applicability.

Reliability hedging required technological uncertainty---managers could not select on skill fit because skill requirements were unknown. This condition is met when technology is genuinely novel, with no organizational history of deployment. It is not met when technology is mature and skill requirements are well-documented. The contrast between Gallifrey (novel robots) and Skaro (mature sorters) illustrates this boundary: reliability-based selection emerged at the novel technology site but not at the mature technology site.

Information shielding was sustainable because temporary workers had limited organizational tenure and restricted access to communication networks. Temps at Gallifrey had been with the organization for weeks or months, not years. They lacked the institutional memory that would help them interpret technological change. They lacked the social networks that transmit informal information about organizational dynamics. These characteristics enabled managers to control what information reached them. Workers with longer tenure, better-developed networks, or greater organizational embeddedness would be harder to shield from automation implications.

The compositional shift emerged because labor was fungible---managers had discretion to assign temps across roles without constraint. Temporary work arrangements provide this fungibility by design. Temps can be moved across positions based on operational needs without the constraints that apply to permanent employees. This fungibility enabled managers to channel reliable workers into robot-adjacent roles. In contexts where workers have stronger claims to specific positions---through union contracts, seniority systems, or simply longer tenure---managers would have less discretion over assignments.

The permanent employee pattern illuminates these boundaries. Permanent employee tenure decreased at the robot facility despite the general organizational preference for retaining experienced staff during transitions. This apparent contradiction is actually consistent with the mechanism: permanent employees had what temps lacked. With longer organizational tenure, they had witnessed prior rounds of technology change. They had context for interpreting what new technology meant. With better access to information networks, they could assess what robot deployment signaled for the facility's trajectory. They talked to each other, to supervisors, to managers. Information shielding that kept temps uninformed was less effective for permanents.

The divergent patterns---tenure increase for temps, tenure decrease for permanents---provide a within-facility test of the mechanism. If the tenure increase reflected facility-wide improvements in working conditions, permanent employees should also have shown increased tenure. They did not. If the increase reflected the inherent attractiveness of robot work, both worker types should have responded similarly. They did not. The divergence is consistent with the mechanism operating through selection and information asymmetry, conditions that applied to temps but not permanents.

\subsection{Implications for Practice}

The practices documented here have implications for managers, policymakers, and workers navigating technological change.

Reliability hedging may produce stable teams capable of adapting to new technology, but it concentrates automation exposure among workers predisposed to stay regardless of working conditions. This could mask problems with automation implementation that would otherwise surface through turnover. Managers should consider whether reliability hedging serves organizational learning or merely suppresses feedback signals.

Information shielding raises ethical considerations. Managers at Gallifrey deliberately withheld displacement implications from workers, enabling reliability hedging to operate without triggering self-selection or flight. This approach may serve short-term operational goals but creates information deficits that workers would likely prefer to avoid. Organizations implementing automation might consider more transparent communication approaches that provide workers with information to make informed choices.

The compositional shift documented here complicates the dominant displacement narrative for policymakers. Aggregate employment statistics may show stable or even increasing employment at firms deploying automation, but this stability could mask compositional changes in who works with automation rather than automation creating attractive jobs. The tenure increase at Gallifrey does not indicate that workers found robot work appealing; it indicates that low-turnover workers were channeled into robot-adjacent roles.

For workers, the findings suggest that automation exposure can depend substantially on organizational selection practices, not only on worker choices or skills. In contexts with high information asymmetry---like temporary work arrangements---workers may be assigned to automation roles without full understanding of the implications.

The account documented here should also prompt reflection on the ethics of information shielding. Managers at Gallifrey chose not to communicate the labor cost implications of automation to workers. This choice served organizational purposes---enabling reliability-based selection without worker self-selection---but it also kept workers from making informed decisions about their labor allocation. Whether such information asymmetry is justifiable depends on one's normative framework, but its existence and consequences deserve scrutiny. Organizations deploying automation make choices about what to tell workers. Those choices affect what workers know and therefore what decisions they can make. The relationship between automation and labor is not just a matter of what technology does to jobs; it is also a matter of what organizations tell workers about what technology does to jobs.

\subsection{Future Research}

The three-layer mechanism developed here opens several directions for future research. First, the mechanism suggests that reliability-based selection should emerge wherever genuinely novel technology meets temporary labor---not only in warehousing but in healthcare, manufacturing, and professional services deploying unfamiliar automation. The key insight is not that robots specifically produce tenure increases, but that any technology creating skill uncertainty in a temporary labor context should trigger reliability-based selection. Examining this pattern across industries would establish the mechanism's generalizability.

Second, research could examine the temporal dynamics of the selection mechanism. As technology matures and skill requirements become understood, does selection shift from reliability toward skill fit? The transition from novel to mature technology may attenuate the selection mechanism over time. Longitudinal data tracking the same facilities across technology lifecycles could examine this possibility.

Third, research could examine heterogeneity in how organizations respond to technological uncertainty. Our qualitative evidence documents reliability-based selection at one organization, but other organizations facing similar uncertainty might respond differently---through training programs, external hiring of workers with adjacent experience, or trial-and-error approaches with high expected turnover. Variation in organizational responses could illuminate when and why reliability-based selection emerges.

Fourth, research could more directly examine worker awareness and its relationship to turnover. Our argument relies on information asymmetry preventing worker self-selection, but we measure awareness only through interviews. Survey measures of automation awareness administered to workers at varying distances from automation deployments could more precisely examine how information shapes behavior.

Fifth, research could examine long-term outcomes for workers selected into automation roles. Our data track tenure---whether workers stay---but not subsequent outcomes like skill development, career progression, or later labor market success. Workers selected for reliability and assigned to automation roles may accumulate valuable skills, or they may be trapped in roles that lead nowhere as technology evolves.

\subsection{Limitations}

This study has limitations characteristic of mixed-methods field research in a single organizational context. The primary quantitative comparison rests on one novel-technology site; the research design proposes a theoretical mechanism and documents consistent patterns rather than establishing causal identification. We infer information asymmetry from qualitative evidence about manager intentions and worker awareness rather than measuring the information that reached each worker directly. And we propose that selection effects dominate treatment effects in producing the tenure pattern, though the observational design cannot fully separate the two.

The data come from fulfillment operations---a setting where temporary labor dominates, performance is individually measurable, and turnover is extremely high. But the core mechanism we document---that technological novelty triggers reliability-based selection when skill requirements are unknowable and the workforce is temporary---should apply wherever three conditions hold: the technology is genuinely novel (eliminating skill-based selection), the workforce consists of temporary or external labor (restricting available selection criteria to reliability), and managers have discretion over worker assignments (enabling compositional shifts). These conditions characterize settings beyond warehousing: healthcare facilities deploying novel diagnostic automation with contract nurses, manufacturing plants introducing collaborative robots staffed by temp agencies, and professional services firms assigning contingent workers to AI-augmented roles. The specific magnitudes will vary, but the selection dynamic should operate wherever these conditions are met.

\subsection{Conclusion}

Workers at a robot facility stayed longer, not shorter. The findings trace this reversal to organizational practices---reliability hedging, information shielding, and the resulting compositional shift---that produced the counterintuitive tenure pattern. The mechanism requires technological uncertainty, labor fungibility, and information asymmetry---conditions present for temporary workers at Gallifrey but not for permanent employees or at mature technology facilities.

This pattern challenges simple narratives about automation and labor. Neither the optimistic story (automation creates good jobs) nor the pessimistic story (automation eliminates jobs) captures what happened. Workers stayed longer not because robot work was good but because reliable workers were channeled into robot-adjacent roles, and reliable workers are workers who do not quit. The tenure increase is real, and its causes are organizational rather than technological.

The compositional shift reframes debates about automation's labor market effects. Aggregate statistics showing stable or increased employment at firms deploying automation are typically interpreted as evidence that automation complements rather than displaces labor. The mechanism documented here demands caution: stable employment can reflect who is selected into automation roles rather than how workers respond to automation. Selection effects can produce patterns that appear worker-driven but partly reflect organizational practices.

The broader lesson concerns the relationship between organizational practices and worker outcomes in automation contexts. Automation does not simply happen to workers; it is implemented through organizational processes that shape who encounters automation and under what conditions. Managers make decisions about who works with new technology. Those decisions are shaped by the information available to them---including, critically, their uncertainty about what the technology requires. When uncertainty is high, managers default to selection criteria that have implications for who stays and who leaves. Reliability hedging and information shielding illustrate how technological novelty changes selection processes in ways that invert expected worker dynamics---at least temporarily, for specific worker populations, under specific organizational conditions.

The findings also raise questions about what happens as technology matures. Reliability hedging is a response to uncertainty; as organizations accumulate experience with new technology, uncertainty decreases and skill-based selection becomes possible. The compositional shift that produces counterintuitive tenure patterns is therefore a transitional phenomenon, concentrated in the early period of novel technology deployment. But as organizations deploy increasingly novel technologies while relying on contingent workforces---gig work, platform labor, temporary staffing---the window of epistemic impossibility may recur with each new technology generation. The conditions that produce reliability hedging are not disappearing; they are proliferating. Understanding when and how organizational selection mediates technology's workforce effects is not a narrow empirical question about one warehouse. It is a central challenge for organizational theory in an era of accelerating technological change.

\bibliographystyle{apalike}
\bibliography{references}

\end{document}
